\section{Họ tiên đề: Sensitivity-n và infidelity}
\label{sec:ch08-ho-tien-de}

\index{Sensitivity-n}%
\index{infidelity}%
Họ này không xếp hạng gì cả, nên nó không vào được ô thứ nhất của vòng lặp mà
bỏ luôn cả vòng. Thay vì lấy $k$ đặc trưng
điểm cao nhất, ta lấy \emph{mọi} tập con cỡ $n$, hoặc lấy mẫu ngẫu nhiên trong
số đó, rồi đòi một đẳng thức đúng trên tất cả. Cách đặt vấn đề này đến từ lý
thuyết chứ không từ thực nghiệm, và nó lặp lại kiểu lập luận mà
chương~\ref{ch:shap} và chương~\ref{ch:attribution-theo-gradient} đã dùng khi
phát biểu các tiên đề của SHAP và của Integrated Gradients.

Ancona và cộng sự phát biểu điều kiện dưới tên Sensitivity-n: một phương pháp
attribution thỏa Sensitivity-n khi tổng các attribution trên một tập con đặc
trưng cỡ $n$ bất kỳ bằng đúng biến thiên của đầu ra do việc bỏ tập con ấy gây
ra~\autocite{msensn}. Với $n$ bằng tổng số đặc trưng, điều kiện này trở thành
\tn{tính đầy đủ}{completeness} mà chương~\ref{ch:attribution-theo-gradient} đã
gặp ở Integrated Gradients; Sensitivity-n là dạng tổng quát của nó cho mọi cỡ tập con.

Yeh và cộng sự viết một điều kiện gần như vậy dưới dạng một sai số cần cực
tiểu hóa. Infidelity của một lời giải thích $\Phi$ là kỳ vọng bình phương độ
lệch giữa biến thiên mà lời giải thích tiên đoán và biến thiên mà mô hình thật
sự cho, lấy trên một phân phối \tn{phép nhiễu}{perturbation}:

\begin{equation}
\label{eq:ch08-infidelity}
\mathrm{INFD}(\Phi, f, x)
= \mathbb{E}_{I \sim \mu_I}
  \left[
    \left(
      I^{\top}\Phi(f,x) - \bigl(f(x) - f(x - I)\bigr)
    \right)^{2}
  \right],
\end{equation}

trong đó $I$ là một biến ngẫu nhiên mà các tác giả mô tả là các phép nhiễu đáng
quan tâm quanh $x$~\autocite{minfid}. Cùng bài báo còn định nghĩa max-sensitivity,
đo lời giải thích thay đổi nhiều nhất bao nhiêu khi đầu vào xê dịch trong một
quả cầu bán kính $r$~\autocite{minfid}. Đại lượng thứ hai này đo
\tn{tính ổn định}{stability}
của lời giải thích, không đo quan hệ của nó với mô hình, nên nó không thuộc
định nghĩa~\ref{def:ch08-chi-so-trung-thuc}.

Họ này mạnh hơn họ xóa dần ở một điểm và yếu đi ở một điểm khác. Nó mạnh hơn vì
điều kiện phải đúng trên nhiều tập con chứ không chỉ trên một thứ tự, nên một
phương pháp không thể qua được bằng cách xếp hạng khéo. Chỗ yếu đi thì tinh hơn: điều kiện được
kiểm giữa lời giải thích và \emph{phản ứng của mô hình trước phép nhiễu}, chứ
không phải giữa lời giải thích và cơ chế mà mô hình thật sự dùng. Nếu phép nhiễu đưa mô hình ra
ngoài phân phối thì $f(x - I)$ trong phương trình~\ref{eq:ch08-infidelity} chính
là đại lượng mà mục trước vừa cho thấy là không đọc được, và một phương pháp
attribution khớp thật khít với nó sẽ đạt infidelity thấp.

Ở đây phân phối $\mu_I$ nhận lại vai trò của độ đo trong
chương~\ref{ch:attribution-tu-nguyen-ly-dau}. Các tác giả mô tả nó là các phép
nhiễu đáng quan tâm và để người dùng chọn~\autocite{minfid}; bài báo không biện
minh một lựa chọn cụ thể nào. Sensitivity-n thì chốt cứng giá trị thay thế bằng
cách đặt các đặc trưng bị bỏ về không~\autocite{msensn}. Vậy dù họ này không đi
qua vòng lặp, nó vẫn phải trả lời câu hỏi thứ hai của vòng lặp, và
mục~\ref{sec:ch08-huan-luyen-lai} vừa cho thấy không có câu trả lời trung lập
nào cho câu hỏi ấy.

Một điều cần nói rõ vì sự vắng mặt của nó cũng là dữ kiện. Bài báo của Ancona và
cộng sự không có câu nào phát biểu giả định mà Sensitivity-n dựa vào để con số
của nó có nghĩa là độ trung thực. Chương này vì thế không gán cho các tác giả
một giả định nào. Điều ghi nhận được là điều kiện ấy được đề xuất như một tính
chất đáng có, và nó chưa từng được đối chiếu với ground truth độc lập.
